%--------------------------------------------------------------
% thesis.tex 
%--------------------------------------------------------------
% Corso di Laurea in Informatica 
% http://if.dsi.unifi.it/
% @Facolt\`a di Scienze Matematiche, Fisiche e Naturali
% @Universit\`a degli Studi di Firenze
%--------------------------------------------------------------
% - template for the main file of Informatica@Unifi Thesis 
% - based on Classic Thesis Style Copyright (C) 2008 
%   Andr\'e Miede http://www.miede.de   
%--------------------------------------------------------------
\documentclass[twoside,openright,titlepage,fleqn,
	headinclude,12pt,a4paper,BCOR5mm,footinclude]{scrbook}
%--------------------------------------------------------------
\newcommand{\myItalianTitle}{Titolo italiano\xspace}
\newcommand{\myEnglishTitle}{Titolo inglese\xspace}
% use the right myDegree option
\newcommand{\myDegree}{Corso di Laurea in Informatica\xspace}
%\newcommand{\myDegree}{
	%Corso di Laurea Specialistica in Scienze e Tecnologie 
	%dell'Informazione\xspace}
\newcommand{\myName}{Nome candidato\xspace}
\newcommand{\myProf}{Relatore\xspace}
\newcommand{\myOtherProf}{Correlatore\xspace}
\newcommand{\mySupervisor}{Nome Cognome\xspace}
\newcommand{\myFaculty}{
	Scuola di Scienze Matematiche, Fisiche e Naturali\xspace}
\newcommand{\myUni}{\protect{
	Universit\`a degli Studi di Firenze}\xspace}
\newcommand{\myLocation}{Firenze\xspace}
\newcommand{\myTime}{Anno Accademico 2014-2015\xspace}
\newcommand{\myVersion}{Version 0.1\xspace}
%--------------------------------------------------------------
\usepackage[italian]{babel}
\usepackage[latin1]{inputenc} 
\usepackage[T1]{fontenc} 
\usepackage[square,numbers]{natbib} 
\usepackage[fleqn]{amsmath}  
\usepackage{ellipsis}
\usepackage{listings}
\usepackage{subfig}
\usepackage{caption}
\usepackage{appendix}
\usepackage{siunitx}
%--------------------------------------------------------------
\usepackage{dia-classicthesis-ldpkg}
%--------------------------------------------------------------
% Options for classicthesis.sty:
% tocaligned eulerchapternumbers drafting linedheaders 
% listsseparated subfig nochapters beramono eulermath parts 
% minionpro pdfspacing
\usepackage[eulerchapternumbers,linedheaders,subfig,beramono,eulermath,
parts]{classicthesis}
%--------------------------------------------------------------
\newlength{\abcd} % for ab..z string length calculation
% how all the floats will be aligned
\newcommand{\myfloatalign}{\centering} 
\setlength{\extrarowheight}{3pt} % increase table row height
\captionsetup{format=hang,font=small}
%--------------------------------------------------------------
% Layout setting
%--------------------------------------------------------------
\usepackage{geometry}
\geometry{
	a4paper,
	ignoremp,
	bindingoffset = 1cm, 
	textwidth     = 13.5cm,
	textheight    = 21.5cm,
	lmargin       = 3.5cm, % left margin
	tmargin       = 4cm    % top margin 
}

\lstset{
  	frame=tb,
	language=Matlab,
  	aboveskip=3mm,
  	belowskip=3mm,
  	showstringspaces=false,
  	columns=flexible,
  	basicstyle={\small\ttfamily},
  	numbers=none,
  	breaklines=true,
  	breakatwhitespace=true,
  	tabsize=3
}
%--------------------------------------------------------------
\begin{document}
\frenchspacing
\raggedbottom
\pagenumbering{roman}
\pagestyle{plain}
%--------------------------------------------------------------
% Frontmatter
%--------------------------------------------------------------
%--------------------------------------------------------------
% titlepage.tex (use thesis.tex as main file)
%--------------------------------------------------------------
\begin{titlepage}
	\begin{center}
   	\large
      \hfill
      \vfill
      \begingroup
         \includegraphics[scale=0.15]{logo/LOGO}\\
%			\spacedallcaps{\myUni} \\
			\myFaculty \\
			\myDegree \\
			\vspace{0.5cm}
         \vspace{0.5cm}
         Tesi di Laurea
      \endgroup
      \vfill
      \begingroup
      	\color{Maroon}\spacedallcaps{\myItalianTitle} \\ $\ $\\
      	\spacedallcaps{\myEnglishTitle} \\
	\bigskip
      \endgroup
      \spacedlowsmallcaps{\myName}
      \vfill
      \vfill
      Relatore: \emph{\myProf}\\
      \vfill
      \vfill
      \myTime
      \vfill
	\end{center}
\end{titlepage}
%--------------------------------------------------------------
% back titlepage
%--------------------------------------------------------------
   \newpage
	\thispagestyle{empty}
	\hfill
	\vfill
	\noindent\myName:
	\textit{\myItalianTitle,}
	\myDegree, \textcopyright\ \myTime
%--------------------------------------------------------------
% back titlepage end
%--------------------------------------------------------------

\pagestyle{scrheadings}
%--------------------------------------------------------------
% Mainmatter
%--------------------------------------------------------------
\pagenumbering{arabic}
% use \cleardoublepage here to avoid problems with pdfbookmark
%\include{intro} % use \myChapter command instead of \chapter
\tableofcontents
\listoffigures
\cleardoublepage
\thispagestyle{empty}
\begin{flushright}
\null\vspace{\stretch {1}}
\emph{"Inserire citazione" \break --- Inserire autore citazione} \vspace{\stretch{2}}\null
\end{flushright}
\cleardoublepage
\chapter{esercizi}

\begin{enumerate}
\medskip
\item 
Scrivere le possibili evoluzioni del programma 
\begin{verbatim}
                       		co X: = X+2 // X: = X+1 oc
\end{verbatim}
assumendo che ciascun assegnamento è realizzato da tre azioni atomiche che
caricano X in un registro (\verb"Load R X"), incrementano il valore del registro (\verb"Add R v") e memorizzano il valore del registro((\verb"Store R X"). Per ciascuna delle esecuzioni risultanti dallÕinterleaving delle azioni atomiche descrivere il contenuto dopo ogni passo della locazione condivisa X e dei registri privati,  $R_1$ del processo che esegue il primo assegnamento ed  $R_2$  per il processo che esegue il secondo assegnamento. Se assuma che il valore iniziale di X sia 50.
\medskip
\item
Si definisca il problema della \emph{barrier synchronization} e si descrivano per sommi capi i differenti approcci alla sua soluzione. Se ne fornisca quindi una soluzione dettagliata utilizzando i semafori.
\medskip
\item
Considerare n api ed un orso che possono avere accesso ad una tazza di miele
inizialmente  vuota e con una capacità di k porzioni. L'orso dorme finchè la
tazza è piena di k-porzioni,  quindi mangia tutto il miele e si rimette a
dormire. Le api riforniscono in continuazione la  tazza con una porzione di
miele finchè non si riempie; l'ape che aggiunge la k-esima porzione sveglia
l'orso. Fornire una soluzione al problema modellando orso ed api come  processi
e utilizzando un monitor per gestire le loro operazioni sulla tazza. Prevedere che le api possano eseguire l'operazione \emph{produce-honey} anche concorrentemente.
 \medskip
 \item  Descrivere le primitive di scambio messaggi send e receive sia sincrone che asincrone ed implementare
 \vspace{-0.5cm}
 \begin{verbatim}
       - synch_send(v:int)
       - send(v:int)
       - receive(x:int)
\end{verbatim}
 \vspace{-0.3cm}
utilizzando le primitive di LINDA.

\end{enumerate}

\begin{thebibliography}{99}

\bibitem{corbellini}
	A. Corbellini, C. Mateos, A. Zunino, D. Godoy, S. Schiaffino - 
	\emph{Persisting Big Data: The NoSQL landscape} - 
	ISISTAN (CONICET-UNCPBA) Research Institute1, UNICEN University, Campus Universitario, Tandil B7001BBO, Argentina - 
	2016

\bibitem{nascita}
	Stefano Bernardi - 
	\emph{An interview with Salvatore Sanfilippo, creator of Redis, working out of Sicily} - 
	\url{http://www.eu-startups.com/2011/01/an-interview-with-salvatore-sanfilippo-creator-of-redis-working-out-of-sicily/}

\bibitem{murmur2flaw}
	\emph{MurMur2 Hash Function Flaw} -
	\url{https://github.com/aappleby/smhasher/wiki/MurmurHash2Flaw}

\bibitem{select}
	\emph{Waiting for Input or Output} - 
	The GNU C Library - 
	\url{https://www.gnu.org/software/libc/manual/html_node/Waiting-for-I_002fO.html}

\bibitem{epoll}
	\emph{epoll(4)} - 
	Linux Man page - 
	\url{https://linux.die.net/man/4/epoll}

\bibitem{aws-latency}
	Steve Newman -
	\emph{Three latency anomalies} - 
	Am I stronger Yet? - 
	\url{http://amistrongeryet.blogspot.it/2010/04/three-latency-anomalies.html}

\bibitem{lua-angry}
	Ansca Inc. -
	\emph{Lua: the lingua franca of iPhone games} -
	Company blog -
	\url{https://web.archive.org/web/20120502071633/http://blog.anscamobile.com/2010/04/lua-the-lingua-franca-of-iphone-games/} -
	Archived by the Internet Web Archive

\bibitem{lua-wow}
	\emph{Lua} -
	WOWWiki on Wikia -
	\url{http://wowwiki.wikia.com/wiki/Lua}

\bibitem{tcp-cork}
	Chris Baus -
	\emph{TCP\_CORK: More than you ever wanted to know} -
	\url{http://baus.net/on-tcp_cork}

\bibitem{tcp-cork-redispy}
	Sorgenti di redispy - connection.py -
	\url{https://github.com/andymccurdy/redis-py/blob/8a186ebd9cecf466624803bdf215030158c7c776/redis/connection.py\#L517}

\bibitem{hyperloglog}
	Philippe Flajolet, Éric Fusy, Olivier Gandouet, e Frédéric Meunier -
	\emph{HyperLogLog: the analysis of a near-optimal cardinality estimation algorithm} -
	2007 Conference on Analysis of Algorithms, AofA 07 -
	\url{http://algo.inria.fr/flajolet/Publications/FlFuGaMe07.pdf}

\bibitem{hyperloglog-plusplus}
	Stefan Heule, Marc Nunkesser e Alexander Hall -
	\emph{HyperLogLog in Practice: Algorithmic Engineering of a State of The Art Cardinality Estimation Algorithm} -
	EDBT/ICDT ’13 March 18 - 
	\url{https://static.googleusercontent.com/media/research.google.com/en//pubs/archive/40671.pdf}

\bibitem{hyperloglog-explain}
	Juan Lopes -
	Answer to question: \emph{How does the HyperLogLog algorithm work?} -
	Stack Overflow -
	\url{https://stackoverflow.com/a/12734343/118015}

\bibitem{fb-mau}
	Josh Cosentine -
	\emph{Facebook now has 2 billion monthly users... and responsibility} -
	TechCrunch -
	\url{https://techcrunch.com/2017/06/27/facebook-2-billion-users/}

\bibitem{lasvegas}
	László Babai -
 	Monte-Carlo algorithms in graph isomorphism testing -
 	Eötvös University, Budapest, Hungary

\bibitem{bloomfilters}
	Burton H. Bloom -
	Space/Time Trade-offs in Hash Coding with Allowable Errors -
	Commun. ACM 13(7): 422-426 -
	1970

\bibitem{golomb}
	Golomb-Rice Coding -
	\url{https://en.wikipedia.org/wiki/Golomb_coding}

\bibitem{golomb-safebrowsing}
	Google Safe Browsing - Compression
	\url{https://developers.google.com/safe-browsing/v4/compression}

\bibitem{bloomnetwork}
	Andrei Broder and Michael Mitzenmacher -
	\emph{Network Applications of Bloom Filters: A Survey} -
	Internet Mathematics Vol. 1, No. 4: 485-509
	\url{http://www.eecs.harvard.edu/~michaelm/postscripts/im2005b.pdf}

\bibitem{bloomproxy}
	L. Fan, P. Cao, J. Almeida, and A. Z. Broder. -
	\emph{Summary Cache: A Scalable Wide-Area Web Cache Sharing Protocol.} -
	IEEE/ACM Transactions on Networking 8:3 (2000), 281—293.

\bibitem{bloomakamai}
	Maggs, Bruce M.; Sitaraman, Ramesh K. -
	\emph{Algorithmic nuggets in content delivery} -
	SIGCOMM Computer Communication Review, New York, NY, USA -
	ACM, 45 (3): 52–66, doi:10.1145/2805789.2805800 -
	July 2015

\bibitem{bloomdoublehash}
	Dillinger, Peter C.; Manolios, Panagiotis - 
	\emph{Bloom Filters in Probabilistic Verification} -
	Proceedings of the 5th International Conference on Formal Methods in Computer-Aided Design, Springer-Verlag, Lecture Notes in Computer Science 3312

\bibitem{bloomscalable}
	Paulo Sergio Almeida, Carlos Baquero, Nuno Preguiça, David Hutchison -
	\emph{Scalable Bloom Filters} -
	Information Processing Letters archive Volume 101 Issue 6, pages 255-261 -
	March 2007 -
	\url{http://gsd.di.uminho.pt/members/cbm/ps/dbloom.pdf}

\bibitem{lemirereduction}
	Daniel Lemire -
	\emph{A fast alternative to the modulo reduction} -
	Daniel Lemire's blog -
	June 2016 -
	\url{https://lemire.me/blog/2016/06/27/a-fast-alternative-to-the-modulo-reduction/}

\end{thebibliography}

%--------------------------------------------------------------
\end{document}
%--------------------------------------------------------------
